\subsection{An incomplete metric with compatible flow data}
\label{sec:flow-compatible-incomplete-metric}

The auxiliary divisorial equation can be incomplete even when its class, divisor and volume come from the prescribed initial data of a geometric smoothing problem.

Let $(X,\omega)$ be compact K\"ahler, with an $\omega$-psh initial potential $\phi_0$ having analytic singularities. The compatible data on a log resolution $\pi:Y\to X$ satisfy
\begin{gather}
 \pi^*\phi_0=\sum_jm_j\log|\sigma_j|^2+\ell_0,\qquad
 \pi^*(\omega+\ddc\phi_0)=\sum_jm_j[D_j]+R_0,
 \label{eq:compatible-initial-resolution}\\
 K_Y=\pi^*K_X+\sum_jb_jD_j,\qquad
 \pi^*\mathrm{d}V_X=\prod_j|\sigma_j|^{2b_j}\mathrm{d}V_Y.
 \label{eq:compatible-volume-resolution}
\end{gather}
Here $D=\sum_jD_j$ is SNC, the defining sections have smooth Hermitian squared norms, $\ell_0$ is smooth, and $R_0$ is smooth closed semipositive. If $\Theta_j$ denotes the curvature of the section norm, put
\[
 \lambda=\max_j\frac{m_j}{1+b_j},\qquad
 \Theta=\sum_j(1+b_j)\Theta_j,\qquad
 \theta=\lambda\Theta+R_0,\qquad
 |\sigma|^2=\prod_j|\sigma_j|^2.
\]
The number $\lambda$ is the inverse exponential-integrability index of $\phi_0$. The metric in the next theorem is the solution of the auxiliary elliptic equation at this background.

\begin{theorem}\label{thm:flow-compatible-incomplete-metric}\index{Metric completeness!compatible initial data}
There are compatible data \eqref{eq:compatible-initial-resolution}--\eqref{eq:compatible-volume-resolution} on smooth projective surfaces, with $\mathrm{d}V_X=\omega^2$, $R_0$ K\"ahler and $D$ a disjoint union of two smooth curves, such that
\begin{equation}\label{eq:compatible-incomplete-equation}
 \langle(\theta+\ddc\rho)^2\rangle
 =\mathrm{e}^{\rho}|\sigma|^{-2}\,\mathrm{d}V_Y
\end{equation}
has a unique full-mass solution in $\mathcal E^1(Y,\theta)$. It has the Lelong numbers of a minimal potential, but every dense Zariski open positive regular locus of its current is incomplete. The unweighted pole is nonintegrable, with smooth strictly positive coefficient across both components of $D$.
\end{theorem}

The auxiliary class in this example is non-nef. Thus the theorem and \cref{thm:nef-incomplete-cusp} answer the two separate completeness assertions negatively. Completeness under simultaneous nefness and these compatibility identities remains open. Neither example identifies its auxiliary metric with the evolving maximal flow.

\Needspace{12\baselineskip}
\begin{lemma}\label{lem:fixed-volume-curve-cusp}
On $\mathbb P^1$, let $z=Z_1/Z_0$, $r=|z|^2$, $\alpha=\ddc\log(1+r)$ and $h=r/(1+r)$, with $\int\alpha=1$. For every $c>2$ there is a rotational $c\alpha$-psh function $u$, smooth off $0$, such that
\begin{equation}\label{eq:fixed-volume-curve-cusp}
 g_1:=c\alpha+\ddc u=\mathrm{e}^{u}h^{-1}c\alpha.
\end{equation}
It has zero Lelong number at $0$, mass $\int g_1=c$, and
\[
 \int_{\mathbb P^1}|u|\alpha+\int_{\mathbb P^1}|u|g_1<\infty.
\]
For $s=-\log r\to\infty$, $t=\log r$ and some $A\in\R$,
\[
 \begin{aligned}
 u&=\log(2/c)-2\log(s+A)+O(\mathrm{e}^{-s}),\\
 u_t&=\frac{2}{s+A}+O(\mathrm{e}^{-s}),\qquad
 u_{tt}=\frac{2}{(s+A)^2}+O(\mathrm{e}^{-s}).
 \end{aligned}
\]
\end{lemma}

\begin{proof}
\emph{Barriers and the radial inverse.}
Set $L=1-\log h=1+\log(1+1/r)$ and $w=-2\log L$. Differentiation gives
\[
 \ddc w=\left(-\frac2L+\frac2{rL^2}\right)\alpha,\qquad
 Q(r):=\frac{c\alpha+\ddc w}{\mathrm{e}^{w}h^{-1}c\alpha}
 =\frac{crL^2-2rL+2}{c(1+r)}.
\]
The function $Q$ is positive, since $L\geq1$ and $c>2$. Its limits are $2/c$ at $0$ and $1-2/c$ at infinity, so $0<q_-\leq Q\leq q_+<\infty$. Choose $a_0\leq\min(0,\log q_-)$ and $b_0\geq\max(0,\log q_+)$. Then $l=w+a_0$ and $H=w+b_0$ satisfy
\[
 c\alpha+\ddc l\geq\mathrm{e}^{l}h^{-1}c\alpha,\qquad
 c\alpha+\ddc H\leq\mathrm{e}^{H}h^{-1}c\alpha,\qquad l\leq w\leq H.
\]

Exhaust $\mathbb P^1\setminus\{0\}$ by disks $|\zeta|<R$ in $\zeta=1/z$, with boundary value $w$. For $q=|\zeta|^2$, the equation becomes
\[
 \Delta u=f(q)\mathrm{e}^{u}-g(q),\qquad
 f(q)=\frac{4c}{1+q},\qquad g(q)=\frac{4c}{(1+q)^2},
\]
where $\Delta$ is the Euclidean real Laplacian.

For $k>0$, the regular radial solution $v=V(q)$ of $(-\Delta+k)v=F(q)$ satisfies
\begin{equation}\label{eq:radial-cusp-volterra}
 V(q)=A_0+\frac14\int_0^q\log(q/t)\bigl(kV(t)-F(t)\bigr)\,\mathrm{d}t.
\end{equation}
The $j$-fold iterate of the integral kernel on $1$ is $q^j/(4^j(j!)^2)$, so successive integration converges uniformly on every finite interval. The coefficient of $A_0$ is $\sum_{j\geq0}(kq/4)^j/(j!)^2>0$. Hence a unique $A_0$ gives any prescribed value at $q=R^2$. Moreover,
\[
 V'(q)=\frac14\int_0^1\bigl(kV(qt)-F(qt)\bigr)\,\mathrm{d}t.
\]
This formula gives smoothness at $q=0$ by induction. The maximum principle for $-\Delta+k$ proves uniqueness and order preservation of the radial inverse.

Choose $k\geq\sup f\mathrm{e}^{H}$ on the fixed disk. Invert $-\Delta+k$ with boundary value $w$ and right side $kv-f\mathrm{e}^{v}+g$, which is nondecreasing on $[l,H]$. The barrier inequalities give $l\leq T(l)$ and $T(H)\leq H$. Iteration from $l$ therefore increases and stays below $H$. The displayed derivative formula gives a uniform first-derivative bound on the fixed interval. Compactness and monotonicity give uniform convergence; passing to \eqref{eq:radial-cusp-volterra} proves the equation, and the averaged derivative formula bootstraps the limit to smoothness.

The same barriers bound the solutions on every fixed reciprocal disk as $R\to\infty$. For an exact nonlinear solution $U(q)$,
\[
 U'(q)=\frac14\int_0^1\bigl(f(qt)\mathrm{e}^{U(qt)}-g(qt)\bigr)\,\mathrm{d}t.
\]
This gives all derivative bounds on compact intervals by induction. A diagonal limit is smooth on the reciprocal plane, including $z=\infty$, and has $u=-2\log(-\log r)+O(1)$ at $0$.

\emph{The asymptotics, mass and energy.}
The radial equation is
\[
 u_{tt}=\frac{c\mathrm{e}^{u}}{1+\mathrm{e}^t}
              -\frac{c\mathrm{e}^t}{(1+\mathrm{e}^t)^2}.
\]
Put $v(t)=u(t)+c\log(1+\mathrm{e}^t)$ and $V(s)=v(-s)$. Then
\[
 V''=k_0(s)\mathrm{e}^{V},\qquad
 k_0(s)=c(1+\mathrm{e}^{-s})^{-(c+1)},\qquad
 V=-2\log s+O(1).
\]
Convexity and this sublinear decay imply $V'<0$ and $V'\to0$ at infinity. The function $J=(V')^2-2k_0\mathrm{e}^{V}$ satisfies $J'=-2k_0'\mathrm{e}^{V}$ and tends to zero. Since $k_0'=O(\mathrm{e}^{-s})$ and $\mathrm{e}^{V}\asymp s^{-2}$,
\[
 J(s)=2\int_s^\infty k_0'(\tau)\mathrm{e}^{V(\tau)}\,\mathrm{d}\tau
       =O(\mathrm{e}^{-s}s^{-2}).
\]
Thus $(\mathrm{e}^{-V/2})'=\sqrt{c/2}(1+O(\mathrm{e}^{-s}))$, and integration gives $\mathrm{e}^{-V/2}=\sqrt{c/2}(s+A)+O(\mathrm{e}^{-s})$. This identity and the differential equation give all three asserted asymptotics.

The local potential $u+c\log(1+r)$ extends subharmonically by \cref{lem:qpsh-removable-analytic}. Its radial slope tends to zero at the puncture, so the extension adds no point mass. The density is asymptotic to
\[
 \frac{\mathrm{i}}{2\pi}\frac{2}{r s^2}
          \,\mathrm{d}z\wedge\mathrm{d}\bar z.
\]
Distributional Stokes gives mass $c$. The energy integral at the cusp is bounded by $C\int^\infty(\log s)/s^2\,\mathrm{d}s$; the integral against the smooth form $\alpha$ is finite as well.
\end{proof}

\Needspace{12\baselineskip}
\begin{proof}[\cref{thm:flow-compatible-incomplete-metric}]
\emph{The initial data and the resolution.}
Take $X=\mathbb P^1_Z\times\mathbb P^1_W$, with normalized Fubini--Study forms $\alpha,\beta$, and put
\[
 \omega=6(\alpha+\beta),\qquad \mathrm{d}V_X=\omega^2=72\alpha\wedge\beta.
\]
Let $H=(Z_1=0)$ and $p=([0:1],[1:0])=(\infty,0)$. Define
\[
 \begin{gathered}
 h_H=\frac{|Z_1|^2}{|Z_0|^2+|Z_1|^2},\qquad
 \eta=\frac{|Z_0|^2}{|Z_0|^2+|Z_1|^2}
              +\frac{|W_1|^2}{|W_0|^2+|W_1|^2},\\
 \phi_0=2\log h_H+2\log\eta.
 \end{gathered}
\]
The zero sets are $H$ and $\{p\}$, which are disjoint. Blow up $p$, write $E$ for the exceptional curve, and set $H_Y=\pi^{-1}H$, $D=H_Y+E$. The projective blowup is smooth, compact, connected and K\"ahler, and is an isomorphism off the initial singular locus.

Use $|\sigma_H|^2=\pi^*h_H$ and $|\sigma_E|^2=\pi^*\eta$. Near $p$, put $x=Z_0/Z_1$, $y=W_1/W_0$. In the chart $x=a$, $y=ab$,
\[
 \pi^*\eta=|a|^2A(a,b),\qquad
 A(a,b)=\frac1{1+|a|^2}+\frac{|b|^2}{1+|ab|^2}>0.
\]
This defines a smooth Hermitian metric on $\mathcal O_Y(E)$ across $E$. The second chart is obtained by interchanging $x,y$, and the metric factors satisfy the transition rule.

The sections $Z_0W_0$, $\sqrt2Z_0W_1$, $Z_1W_1$ of $\mathcal O(1,1)$ have summed squared norm $\eta$ and base ideal $(x,\sqrt2xy,y)=(x,y)$. Dividing by $\sigma_E$ gives a morphism $F:Y\to\mathbb P^2$ with
\begin{equation}\label{eq:compatible-curvature-map}
 F^*\gamma=\pi^*(\alpha+\beta)-\Theta_E,\qquad
 \Theta_H=\pi^*\alpha,
\end{equation}
where $\gamma$ is normalized Fubini--Study. In the first chart $F=[1:\sqrt2ab:b]$, so $F|_E$ is an isomorphism onto a line.

Near $p$,
\[
 \eta=\frac{|x|^2+|y|^2+2|xy|^2}{(1+|x|^2)(1+|y|^2)}.
\]
Thus $\phi_0$ has analytic singularities; in the convention $\frac1{2m}\log\sum|f_j|^2$, the coefficient $2$ is obtained by taking the fourth power of the sum of squares and expanding it as a sum of squared holomorphic monomials. The exact resolved identities are
\begin{align}\label{eq:compatible-kahler-residual}
 \pi^*\phi_0&=2\log|\sigma_H|^2+2\log|\sigma_E|^2,\\
 R_0&=\pi^*\omega-2\Theta_H-2\Theta_E
       =2\pi^*\alpha+4\pi^*\beta+2F^*\gamma.\notag
\end{align}
The first two terms of $R_0$ are positive away from $E$. On $E$, their kernel is $TE$, and $F^*\gamma$ is positive on $TE$. Hence $R_0$ is K\"ahler. The resolved current is $2[H_Y]+2[E]+R_0$, which proves $\omega$-psh positivity away from $p$; the bounded-above extension across $p$ gives it everywhere. We have $m_H=m_E=2$ and $\ell_0=0$.

The blowup Jacobian is $a$, so $b_E=1$, $b_H=0$. Direct exponential integrability requires $q<1/2$ at $H$ and $q<1$ at $p$. Consequently $\lambda=2=\max\{2,1\}$, and
\begin{equation}\label{eq:compatible-auxiliary-background}
 \Theta=\Theta_H+2\Theta_E,\qquad
 \theta=2\Theta+R_0=\pi^*\omega+2\Theta_E.
\end{equation}
Since $\Theta_E|_E=-\omega_{\mathrm{FS}}$, the class has intersection $-2$ with $E$ and is non-nef. It is big: the current associated to $3\log|\sigma_E|^2$ is
\[
 (\pi^*\omega-\Theta_E)+3[E],\qquad
 \pi^*\omega-\Theta_E=5\pi^*(\alpha+\beta)+F^*\gamma>0.
\]

\emph{The exact prescribed volume.}
Define $\mathrm{d}V_Y=\pi^*\omega^2/|\sigma_E|^2$ off $E$. In the chart above,
\[
 \mathrm{d}V_Y=\frac{72\,\mathrm{d}A_a\mathrm{d}A_b}
 {\pi^2A(a,b)(1+|a|^2)^2(1+|ab|^2)^2}.
\]
The cancellation of the Jacobian $|a|^2$ is exact. The second chart gives the extension on the rest of $E$, and the extensions agree on overlaps. Thus the volume is globally smooth and strictly positive, with $\pi^*\mathrm{d}V_X=|\sigma_E|^2\mathrm{d}V_Y$ as required. Since $H_Y$ and $E$ are disjoint, $|\sigma_H\sigma_E|^{-2}\mathrm{d}V_Y$ has a nonintegrable ordinary pole with smooth positive coefficient at each component.

\emph{The solution and the minimal reference.}
Take $u$ from \cref{lem:fixed-volume-curve-cusp} with $c=6$, pulled to the first factor of $X$, and put $g_X=(6\alpha+\ddc u)+6\beta$. It is smooth and strictly positive near $p$, and
\[
 g_X^2=\mathrm{e}^{u}h_H^{-1}\omega^2.
\]
Set $P=2\log|\sigma_E|^2$ and $\rho=P+\pi^*u$. Then
\[
 \theta+\ddc P=2[E]+\pi^*\omega,\qquad
 \theta+\ddc\rho=2[E]+\pi^*g_X.
\]
Off $D$ the equation follows from the exact identity
\[
 \mathrm{e}^{\rho}|\sigma_H\sigma_E|^{-2}\mathrm{d}V_Y
 =\mathrm{e}^{\pi^*u}|\sigma_H|^{-2}|\sigma_E|^2\mathrm{d}V_Y
 =\pi^*g_X^2.
\]
Both measures extend by zero on $D$, proving the global non-pluripolar equation.

For any $\theta$-psh potential $\psi$, the function $\psi-P$ descends off $E$ to an $\omega$-psh function on $X\setminus\{p\}$. By \cref{lem:isolated-puncture-extension}, it extends and is bounded above on $X$. Thus $\psi\leq P+C_\psi$ off $E$, and \cref{lem:qpsh-almost-everywhere-order} gives the inequality everywhere. Since $P$ itself is $\theta$-psh, it has minimal singularities.

Its current has non-pluripolar mass $\int_X\omega^2=72$, the volume of the class. In contrast, the ordinary cup-product square is $72+4E^2=68$. The cusp has mass $6$ and no atom, so the current at $\rho$ also has mass $\int_Xg_X^2=72$. Thus $\rho$ has full mass. Near $E$, $\pi^*u$ is smooth, giving the Lelong numbers of $P$; on $H_Y$ the log-log cusp has zero Lelong number. The actual solution is not of ordinary minimal type, because $u$ is unbounded below at $H$.

\emph{Relative energy and uniqueness.}
Choose $M$ with $u-M\leq0$ and use the minimal reference $P+M$. Its canonical approximants are
\[
 \rho_k=P+M+\pi^*u_k,\qquad u_k=\max(u-M,-k).
\]
Their relative difference off $E$ is $\pi^*u_k$. On the curve let $b(t)=6\log(1+\mathrm{e}^t)$ and $V(t)=b(t)+u(t)$. For large $k$ there is a unique cusp level $t=T_k$ with $u(T_k)-M=-k$. If $S_k=-T_k$, the asymptotics give $S_k\asymp\mathrm{e}^{k/2}$.

The truncated curve primitive is $b+u_k=\max(V-M,b-k)$. Below $T_k$ its mass is $b'(T_k)$; at $T_k$ the derivative jumps upward by $u_t(T_k)$. Smooth convex approximation and bounded Bedford--Taylor continuity therefore give total cutoff mass $V'(T_k)=O(S_k^{-1})$. On that region $u_k=-k$, so its absolute energy contribution is
\[
 kV'(T_k)=O(k\mathrm{e}^{-k/2})\longrightarrow0.
\]
Outside it the original density and potential are unchanged, and $\int|u|g_1$ is finite. Consequently $\int|u_k|g_{1,k}$ has a finite limit, where $g_{1,k}=6\alpha+\ddc u_k$. Also $\int|u_k|\alpha$ converges by domination.

The mixed measures for the product current and its reference are
\[
 \omega^2=72\alpha\wedge\beta,\quad
 (g_{1,k}+6\beta)\wedge\omega
     =6g_{1,k}\wedge\beta+36\alpha\wedge\beta,\quad
 (g_{1,k}+6\beta)^2=12g_{1,k}\wedge\beta.
\]
Their pullbacks are precisely the mixed non-pluripolar measures upstairs, ignoring $E$. All three canonical relative energy integrals therefore have finite limits. This proves $\rho\in\mathcal E^1(Y,\theta)$. Replacing $P+M$ by the normalized minimal reference changes it by a bounded amount and preserves this finite-energy class.

Uniqueness follows by the same full-mass comparison argument as in \cref{thm:incomplete-divisorial-metric}. The unweighted pole measure has infinite mass, but the two exponential densities being compared have mass $72$. Comparison on the strict sublevel sets forces equality almost everywhere, and then equality of the canonical quasi-psh representatives.

\emph{The missing metric boundary.}
The positive smooth locus of the current is $Y\setminus(H_Y\cup E)$. The cusp coefficient diverges at $H_Y$, and the nonzero current $2[E]$ excludes smoothness at $E$. The pullback metric part is also degenerate along $TE$.

For any dense Zariski open positive regular locus $U$, choose a point of $E$ avoiding every other component of its analytic complement, as in \cref{thm:nef-incomplete-cusp}. In a small chart the complement is exactly $E$. A transverse path $a=\xi>0$, $b=b_0$ lies in $U$ for small $\xi$ and approaches that point. Since $g_X$ is smooth near $p$ and $\pi$ is smooth, the coefficients of $\pi^*g_X$ are bounded in the chart. The path has bounded speed and finite remaining length. Its tails give an intrinsic Cauchy sequence with no limit in $U$, since its only ambient limit is the removed point and a smooth positive Riemannian metric induces the manifold topology. Hence every such $U$ is incomplete.
\end{proof}
